\documentclass[11pt]{article}

\usepackage{amsmath, amssymb, graphicx, hyperref}
\usepackage{geometry}
\geometry{margin=1in}

\title{Derivative-Aware Ray Stepping with Deadband Filtering:\\
A Sparse Control Strategy and the Perceptual Curvature Threshold Hypothesis}

\author{Anonymous Author(s)}
\date{}

\begin{document}

\maketitle

\begin{abstract}
We introduce a deadband-filtered derivative-aware stepping method for curved-ray integration and demonstrate that it eliminates the majority of control adjustments while preserving bitwise-identical image outputs. We propose the \textit{Perceptual Curvature Threshold Hypothesis}, which posits that curvature variations below a finite threshold do not contribute to observable ray trajectory differences under discrete sampling. We connect this threshold to geodesic deviation limits in General Relativity and to Nyquist-like sampling constraints. Experimental results show $10^4$--$10^5$ reduction in applied adjustments with no change in output, suggesting that sparse, thresholded control may be a fundamental principle for efficient curved-ray transport.
\end{abstract}

\section{Introduction}

Curved-ray transport appears in gradient-index (GRIN) optics, atmospheric refraction, and metric-based rendering inspired by General Relativity (GR). Numerical integration of such systems often employs derivative-aware step control to adapt to local curvature.

However, these methods frequently exhibit excessive sensitivity, producing large volumes of micro-adjustments that do not influence the final rendered image.

This work investigates whether a thresholded response can preserve correctness while dramatically reducing control activity.

\section{Method}

\subsection{Derivative-Aware Step Scaling}

Given a local curvature signal $\nabla \mathcal{F}$, a raw scale factor is computed:

\begin{equation}
s_{\text{raw}} = f(\nabla \mathcal{F})
\end{equation}

\subsection{Deadband Filtering}

We introduce a deadband threshold $\epsilon$:

\begin{equation}
s =
\begin{cases}
1.0 & \text{if } |s_{\text{raw}} - 1.0| < \epsilon \\
s_{\text{raw}} & \text{otherwise}
\end{cases}
\end{equation}

This suppresses sub-threshold adjustments.

\section{Experimental Setup}

Fixtures:
\begin{itemize}
\item \texttt{curved\_minimal}
\item \texttt{curved\_minimal\_backdrop}
\end{itemize}

Metrics:
\begin{itemize}
\item Candidate vs applied adjustments
\item Runtime
\item Output image hash
\end{itemize}

\section{Results}

\subsection{Adjustment Reduction}

\begin{itemize}
\item Candidate adjustments: $\sim 10^6$
\item Applied adjustments: $\sim 10^2$
\end{itemize}

\subsection{Output Fidelity}

All runs produced identical image hashes, indicating no observable deviation.

\subsection{Runtime}

Runtime remained approximately constant, indicating that adjustment logic is not the dominant computational cost.

\section{Perceptual Curvature Threshold Hypothesis}

We propose that a finite curvature threshold $\kappa_p$ exists such that:

\begin{equation}
\kappa < \kappa_p \Rightarrow \Delta x_{\text{ray}} < \Delta x_{\text{pixel}}
\end{equation}

where $\Delta x_{\text{ray}}$ is deviation in ray trajectory and $\Delta x_{\text{pixel}}$ is the sampling resolution.

Thus, curvature below $\kappa_p$ is observationally irrelevant.

\section{Connection to Geodesic Deviation}

In GR, geodesic deviation is governed by:

\begin{equation}
\frac{D^2 \xi^\mu}{d\tau^2} = - R^\mu_{\ \nu\alpha\beta} u^\nu \xi^\alpha u^\beta
\end{equation}

Small curvature implies small deviation acceleration. Under discrete integration and finite resolution, these deviations may fall below sampling detectability.

\section{Nyquist-Like Interpretation}

We interpret curvature as a signal with frequency content. If curvature variation occurs below the effective sampling frequency:

\begin{equation}
\kappa < \kappa_{\text{Nyquist}} \Rightarrow \text{undetectable variation}
\end{equation}

This suggests a sampling-theoretic limit on meaningful curvature response.

\section{Discussion}

The deadband acts as a filter enforcing:

\begin{itemize}
\item Sparse control
\item Noise suppression
\item Stability
\end{itemize}

This aligns with principles from signal processing and control theory.

\section{Future Work}

\begin{itemize}
\item Hysteresis-based temporal stability
\item Adaptive threshold estimation
\item Extension to metric-based geodesic solvers
\end{itemize}

\section{Conclusion}

We demonstrate that most derivative-driven adjustments in curved-ray transport are unnecessary under discrete sampling. A simple deadband filter preserves correctness while enabling sparse control, suggesting a fundamental threshold governing observable curvature.

\end{document}